% !TEX root = ./main.tex
\chapter*{Sommario}
\addcontentsline{toc}{chapter}{Sommario}

SOULFRAME nasce dall'idea di costruire un interlocutore embodied che non si limiti a collocare un chatbot in una scena 3D, ma mantenga nel tempo un'identità riconoscibile per aspetto, voce e memoria. Nel prototipo l'utente può creare o ricostruire il proprio avatar tramite Avaturn, associargli una voce a partire da un breve campione e avviare una conversazione vocale dentro un client Unity.

La tesi descrive l'integrazione effettivamente realizzata tra frontend Unity e servizi backend per trascrizione, gestione conversazionale e sintesi vocale. Il profilo avatar comprende anche una memoria persistente per-avatar, arricchibile con note, immagini e PDF, così da mantenere continuità tra sessioni diverse. Il blocco conversazionale recupera il contesto utile da queste sorgenti e lo usa per produrre risposte più aderenti al profilo e alla storia dell'avatar.

Durante lo sviluppo sono emersi problemi concreti di latenza, compatibilità tra ambienti e tenuta del playback audio, soprattutto nel passaggio dal locale al deploy remoto. Per questo il lavoro non si è limitato alla pipeline conversazionale, ma ha incluso anche setup, deploy e gestione dei servizi. Il risultato è un prototipo funzionante, utilizzabile sia in locale sia su server, che mostra con chiarezza quali aspetti rendono più credibile un interlocutore embodied e quali limiti tecnici restano ancora aperti.
